\section{Results}
\label{sec:results}

This section presents the experimental evaluation of the proposed data-driven VSM/VSD framework. The results are organized into four complementary parts. First, the framework is evaluated under controlled conditions using the synthetic manufacturing event log, allowing quantitative analysis of the cumulative framework configurations through offline Monte Carlo simulation. Second, the analytical components are assessed on three heterogeneous real-world event logs for classification and structural description across domains. Third, the scenario comparison quantifies the incremental operational improvements obtained as additional analytical capabilities are progressively incorporated into the framework. Finally, statistical analyses determine whether the observed operational improvements exceed the variability introduced by stochastic simulation.

\subsection{Experimental Results on the Synthetic Manufacturing Event Log}
\label{subsec:synthetic_results}

The proposed data-driven VSM/VSD framework was first evaluated using a
synthetic manufacturing event log designed to represent controlled production
flows, waiting times, rework activities, quality events, and resource
interactions under controlled experimental conditions. Six analytical
configurations were considered, ranging from a baseline representation of
the process to the complete framework integrating PM,
automated VSM, ML, offline MCS, and CNA. The main results are
summarized in Figure~\ref{fig:synthetic_results}.

\begin{figure*}[t]
\centering
\includegraphics[width=\textwidth]{paper_synthetic_results_final_dashboard.png}
\caption{Synthetic manufacturing simulation results. (A) G5 relative KPI changes versus G0, using the values in Table~\ref{tab:synthetic_improvements}. (B) Cumulative lead-time, throughput, and WIP changes from Table~\ref{tab:ablation_results}. All changes concern simulated scenarios, not physical factory measurements.}
\label{fig:synthetic_results}
\end{figure*}

Figure~\ref{fig:synthetic_results}A summarizes the simulated G5 changes relative to G0. As reported in Table~\ref{tab:synthetic_improvements}, lead time, WIP, defect rate, and rework rate decrease, while throughput increases. These values describe the selected synthetic operating conditions and intervention configurations.

\begin{table*}[t]
\caption{
Model responses under the predefined G5 factor set relative to G0. Workload and throughput quantities follow the proxy conventions of the scenario sampler.
}
\label{tab:synthetic_improvements}
\centering
\begin{tabular}{lc}
\toprule
\textbf{Metric} & \textbf{Improvement (\%)} \\
\midrule
Lead Time Reduction & 32.40 \\
Work-in-Process Reduction & 45.09 \\
Defect Rate Reduction & 27.63 \\
Rework Rate Reduction & 45.52 \\
Throughput Increase & 71.73 \\
\bottomrule
\end{tabular}
\end{table*}

Figure~\ref{fig:synthetic_results}B shows progressively larger simulated changes across G0--G5. The comparison is cumulative: analytical capabilities support intervention choices whose operational effects are evaluated in the simulator. Adding an analytical module does not by itself alter physical process behavior, and the design does not isolate independent causal effects of individual modules.

The synthetic setting enables controlled comparison, but it cannot reproduce every source of variability in industrial production. The exact magnitudes therefore remain scenario specific. Manufacturing deployment would require calibration and independent checking against plant observations, followed by measurements of implemented interventions under documented operating conditions.

\subsubsection{Code-Level Audit of the Synthetic Classification Task}
\label{subsec:synthetic_code_audit}

The delivered synthetic data contain 8,532 events from 1,200 cases and eight observed activities. The flag \texttt{rework\_flag} equals the indicator of the activity \texttt{Rework} in every row. In the chronological test partition, 21 of 1,280 rows are positive. A rule learned only from training rows that predicts each activity's majority label attains F1 = 1.0000, as do an activity-only Random Forest and a Random Forest using training-only centralities. Their confusion matrix is TN = 1,259, FP = 0, FN = 0, TP = 21.

\begin{table*}[t]
\caption{Executed synthetic rework audit. All rows use the same 70/15/15 temporal row split; graph features and activity rules are fitted on training rows only. These are diagnostic-label results, not prospective forecasts.}
\label{tab:synthetic_audit}
\centering\small
\begin{tabular}{lr}
\toprule
Condition & Test F1\\
\midrule
Training-only activity-majority rule & 1.0000\\
Random Forest, operational variables & 0.2500\\
Random Forest, training-only network descriptors & 1.0000\\
Random Forest, operational and training-only network descriptors & 1.0000\\
Random Forest, activity one-hot representation & 1.0000\\
\bottomrule
\end{tabular}
\end{table*}

The four centralities give distinct signatures to all eight activities. Although the original synthetic script loads full-log centralities before splitting, reconstructing the graph using training rows alone produces numerically identical values in this dataset. The audit therefore identifies activity/label reconstruction as a sufficient explanation of its perfect rework F1; it does not attribute the score to a measured difference between full-log and training-only centralities. This experiment does not establish the BPI2017 mechanism by itself; the separate real-log audit establishes the association in its prepared rows.

The synthetic delay flag also equals \texttt{waiting\_time\_min > 15} in every row, and the original synthetic model includes that same waiting variable. Its near-perfect delay score consequently reconstructs the generating rule. This is not evidence of anticipating delays before the relevant information exists. The original row split shares two cases between training and validation and three between validation and test; training and test share no cases. Two pre-test events finish at or after the test start. A prospective design would need an explicit observation cutoff and case-boundary policy rather than relying on row ordering alone.

\subsection{Results on Real-World Event Logs}
\label{subsec:real_results}

To evaluate the generalizability of the proposed framework beyond controlled
experimentation, the methodology was applied to three heterogeneous
real-world event logs exhibiting different operational characteristics,
process complexity, and organizational contexts: BPI2017, BPI2019, and the
Data Quality Issues (DQI) dataset. Collectively, these datasets provide a
challenging evaluation environment that includes large-scale business
processes, procurement workflows, and manufacturing operations affected by
quality-related issues.

\begin{figure*}[t]
\centering
\includegraphics[width=\textwidth]{fig_real_datasets_experimental_results_dashboard}
\caption{Analytical results on BPI2017, BPI2019, and DQI. (A) Activity-network density. (B) Random Forest F1-scores with and without CNA features, using Table~\ref{tab:real_ml}; NA indicates a single-class test target. (C) Activities with the greatest reported betweenness centrality. Synthetic simulation results are presented separately. These classification results concern the event-log/proxy tasks described in the text.}
\label{fig:real_dashboard}
\end{figure*}

Table~\ref{tab:dataset_characteristics} summarizes the principal
characteristics of the evaluated event logs. BPI2019 represents the largest
dataset in terms of cases, activities, and organizational resources,
whereas DQI corresponds to a substantially smaller manufacturing process
characterized by pronounced quality-related issues. The considerable
differences in scale, process structure, and operational behavior make
these datasets appropriate for evaluating the robustness and
generalizability of the proposed framework.

\begin{table*}[t]
\caption{
Descriptive characteristics of the processed real-world event logs used for cross-domain analytical assessment. Counts refer to the analyzed subsets, not necessarily the complete public repositories.
}
\label{tab:dataset_characteristics}
\centering
\small
\begin{tabular}{lrrrrr}
\toprule
\textbf{Dataset} &
\textbf{Events} &
\textbf{Cases} &
\textbf{Activities} &
\textbf{Resources} &
\textbf{Observation Hours} \\
\midrule
BPI 2017 & 250,000 & 31,469 & 26 & 147 & 9,530.57 \\
BPI 2019 & 250,000 & 155,771 & 39 & 563 & 622,190.53 \\
Data Quality Issues & 8,607 & 271 & 21 & 15 & 5,977.72 \\
\bottomrule
\end{tabular}
\end{table*}

The descriptive process indicators reveal substantial differences in process
behavior across the evaluated datasets. The processed BPI files have reported processing-time fractions of approximately 28.1\% and 21.2\%. These fractions cannot be interpreted as validated value-added proportions because the preparation and timestamp lifecycle semantics differ by dataset. The BPI2019 start-time range includes 26 January 1948 through 18 January 2019; the 622,190.53-hour recorded span is therefore reported as a data characteristic, not a validated operating horizon for industrial throughput inference.

In contrast, the DQI dataset presents a Value-Added Ratio close to
100\%. This value should not be interpreted as evidence of near-perfect Lean
performance. Rather, in the supplied adapter the next-event gaps sum to case elapsed time, making the processing-time fraction approximately one by construction. The previous-event gaps are also recorded as waiting proxies, so these derived intervals must not be treated as independent measurements of value-adding service and queue time. This interpretation is particularly important
because the same dataset simultaneously exhibits a rework rate of
74.8\%, revealing substantial process instability that cannot be inferred
from conventional time-based Value Stream Mapping indicators alone.
Accordingly, this dataset illustrates the importance of complementing
traditional VSM analysis with quality-related, predictive, and structural
information.

Complex Network Analysis was subsequently applied to characterize structural
relationships among process activities. The resulting activity-network
properties are summarized in
Table~\ref{tab:network_metrics}.

\begin{table*}[t]
\caption{
Structural properties of the activity networks derived from the evaluated
real-world event logs. Reported activities correspond to the nodes with the
highest structural relevance according to betweenness centrality.
}
\label{tab:network_metrics}
\centering
\small
\begin{tabularx}{\textwidth}{lrrrrX}
\toprule
\textbf{Dataset} &
\textbf{Nodes} &
\textbf{Edges} &
\textbf{Density} &
\textbf{WCC} &
\textbf{Structurally Relevant Activity} \\
\midrule

BPI 2017 & 26 & 380 & 0.5846 & 1 & \texttt{W\_Call incomplete files} \\
BPI 2019 & 39 & 527 & 0.3556 & 1 & \texttt{Record Goods Receipt} \\
Data Quality Issues & 21 & 102 & 0.2429 & 1 & \texttt{/vgr/pick\_up\_and\_transport} \\
\bottomrule
\end{tabularx}
\end{table*}

As illustrated in Figure~\ref{fig:real_dashboard}A, BPI2017 exhibits the
highest activity-network density (0.5846), indicating comparatively strong
connectivity among process activities. BPI2019 presents an intermediate
density (0.3556), whereas DQI shows the sparsest network structure
(0.2429). The structurally relevant activities identified through
betweenness centrality analysis occupy strategically important positions
within the observed activity-flow networks and provide complementary
structural information beyond conventional flow-oriented Value Stream
Mapping indicators.

To evaluate the predictive contribution of structural information,
Machine Learning models were trained using two matched experimental
configurations: one employing only operational process features and another
augmenting the same feature set with Complex Network Analysis descriptors.
Following the temporal protocol described in
Section~\ref{subsec:training_protocol}, proxy-label thresholds were fitted on the training subset, and the stated
predictor exclusions were applied before evaluation on the temporally
held-out testing partition.

Table~\ref{tab:real_ml} summarizes the resulting predictive performance.

\begin{table*}[t]
\centering
\caption{
Temporally evaluated event-log classification on the real-world event logs.
F1-scores compare models using operational features only (No CNA) against
the same models augmented with CNA features.
}
\label{tab:real_ml}
\resizebox{\textwidth}{!}{%
\begin{tabular}{lllllll}
\toprule
Dataset & Target & Model & F1 (No CNA) & F1 (CNA) & $\Delta$F1 & Status \\
\midrule
BPI2017 & Delay & Random Forest & N/A & N/A & N/A & Single-class target \\
BPI2017 & Delay & XGBoost & N/A & N/A & N/A & Single-class target \\
BPI2017 & Bottleneck & Random Forest & N/A & N/A & N/A & Single-class target \\
BPI2017 & Bottleneck & XGBoost & N/A & N/A & N/A & Single-class target \\
BPI2017 & Rework & Random Forest & 0.2233 & 1.0000 & 0.7767 & Evaluable \\
BPI2017 & Rework & XGBoost & 0.0000 & 1.0000 & 1.0000 & Evaluable \\
BPI2019 & Delay & Random Forest & N/A & N/A & N/A & Single-class target \\
BPI2019 & Delay & XGBoost & N/A & N/A & N/A & Single-class target \\
BPI2019 & Bottleneck & Random Forest & N/A & N/A & N/A & Single-class target \\
BPI2019 & Bottleneck & XGBoost & N/A & N/A & N/A & Single-class target \\
BPI2019 & Rework & Random Forest & 0.4665 & 0.9608 & 0.4943 & Evaluable \\
BPI2019 & Rework & XGBoost & 0.0062 & 0.7453 & 0.7391 & Evaluable \\
Data Quality Issues & Delay & Random Forest & 0.1560 & 0.2989 & 0.1429 & Evaluable \\
Data Quality Issues & Delay & XGBoost & 0.0000 & 0.2210 & 0.2210 & Evaluable \\
Data Quality Issues & Bottleneck & Random Forest & 0.6421 & 0.9951 & 0.3530 & Evaluable \\
Data Quality Issues & Bottleneck & XGBoost & 0.0000 & 0.9975 & 0.9975 & Evaluable \\
Data Quality Issues & Rework & Random Forest & 0.9352 & 0.9463 & 0.0111 & Evaluable \\
Data Quality Issues & Rework & XGBoost & 0.9412 & 0.9508 & 0.0096 & Evaluable \\
\bottomrule
\end{tabular}%
}
\end{table*}

The table reports the corrected-code execution with mixed-format UTC timestamp parsing and stable temporal row ordering. The original parser silently rejected 467 otherwise valid BPI2017 rows and six otherwise valid DQI rows; all rows are retained by the corrected parser. Stable ordering also makes the tie convention explicit. Original-code reproductions are retained separately in the audit materials. For BPI2017 they reproduce the reported RF change from 0.2243 to 1.0000; the corrected run yields 0.2233 to 1.0000.

The BPI delay and operational-bottleneck test tasks remain single class under the stated proxy definitions and are reported as unavailable. In particular, waiting-time proxies are constant zero in both prepared BPI files, making the upper-percentile waiting rule uninformative. This is a limitation of the prepared representation rather than evidence that no real delays or bottlenecks occurred.

In BPI2017, all positive prepared rework labels occur at \texttt{O\_Returned}, and every occurrence of that activity is positive. On the corrected 37,500-event test set, an activity-majority rule fitted only on training rows attains F1 = 1.0000, as do the activity-one-hot and training-only network classifiers. The perfect score is thus explained by a deterministic activity--label association in the prepared data; a graph descriptor can serve as an activity signature. It is not evidence that network analysis anticipates future rework or outperforms an activity-only diagnostic baseline.

In BPI2019, positive labels identify a subset of change-related activities. The training-derived activity-majority and activity-only classifiers also attain F1 = 1.0000 on the audited test set. The network-augmented classifiers therefore cannot establish added predictive value beyond the available activity identity. Differences from operational-only models quantify a representation change, not a uniquely structural or causal improvement.

The DQI label represents repeated within-case activity occurrences in the supplied preparation. Its current-event timing and repetition semantics differ from the activity-defined BPI labels. The corrected DQI F1 values are retained in Table~\ref{tab:real_ml}, with activity-baseline results in Table~\ref{tab:real_activity_audit}. Small differences from the original manuscript are documented in the reproduction materials rather than hidden by retaining obsolete rounded values.

\begin{table*}[t]
\caption{Executed activity-baseline comparison for the rework/repetition label. Real-log rows use mixed-format timestamp parsing, stable temporal ordering, and balanced Random Forest class weights; the synthetic RF is unweighted, matching its original classifier. All rules and graph descriptors use training rows only.}
\label{tab:real_activity_audit}
\centering\small
\begin{tabular}{lrrrrrr}
\toprule
Dataset & Test events & Positives & Operational RF & + CNA RF & Activity rule & Activity RF\\
\midrule
Synthetic & 1,280 & 21 & 0.2500 & 1.0000 & 1.0000 & 1.0000 \\
BPI2017 & 37,500 & 846 & 0.2233 & 1.0000 & 1.0000 & 1.0000 \\
BPI2019 & 37,500 & 1,276 & 0.4665 & 0.9608 & 1.0000 & 1.0000 \\
DQI & 1,292 & 948 & 0.9352 & 0.9463 & 0.8464 & 0.4817 \\
\bottomrule
\end{tabular}
\end{table*}

The corrected row split shares seven BPI2017 cases and 2,424 BPI2019 cases between training and test; DQI shares none between those partitions. These counts make the evaluation explicitly event-held-out rather than a claim of unseen-case generalization. The metadata and saved predictions identify the observation windows and class counts.

The real-log code constructs its activity graph exclusively from training rows, and unseen activities receive zero centralities. That control is implemented and is sufficient to rule out direct construction of those graph features from held-out rows in this script. It does not establish a future-event forecasting design: current-event durations can depend on later completion information, and cases overlap the row partitions. The revised interpretation is diagnostic classification of the defined prepared labels.

\subsection{Cumulative Scenario Comparison}
\label{subsec:ablation_results}

The scenario comparison evaluates the predefined factor sets G0--G5. Figure~\ref{fig:ablation_improvement} and Table~\ref{tab:ablation_results} summarize the resulting model indicators. These are cumulative scenario assumptions rather than observed module-addition effects.

\begin{figure*}[t]
\centering
\includegraphics[width=\textwidth]{fig_6_ablation_improvement}
\caption{Simulated relative KPI changes across the cumulative configurations. Values follow Table~\ref{tab:ablation_results}; MCS denotes offline Monte Carlo simulation, and the G4 throughput increase is 46.6\%. These are synthetic scenario outcomes rather than measurements after physical implementation.}
\label{fig:ablation_improvement}
\end{figure*}

\begin{table*}[t]
\centering
\caption{
Relative KPI changes across cumulative configurations in the synthetic manufacturing simulation.
}
\label{tab:ablation_results}
\begin{tabular}{lrrr}
\toprule
\textbf{Configuration} &
\textbf{LT reduction} &
\textbf{TH gain} &
\textbf{WIP reduction} \\
& \textbf{(\%)} & \textbf{(\%)} & \textbf{(\%)} \\
\midrule
G0 -- Baseline & 0.0 & 0.0 & 0.0 \\
G1 -- Manual VSM & 4.8 & 8.6 & 8.2 \\
G2 -- PM-VSM & 9.1 & 17.4 & 15.2 \\
G3 -- PM-VSM + XAI & 14.5 & 27.9 & 21.8 \\
G4 -- PM-VSM + XAI + MCS & 22.8 & 46.6 & 33.1 \\
G5 -- Complete framework & 32.4 & 71.7 & 45.1 \\
\bottomrule
\end{tabular}
\end{table*}

The G5 factor set produces the largest favorable changes among the six scenarios. Relative to G0, sampled lead time is 32.40\% lower, the workload proxy is 45.09\% lower, station defect proportion is 27.63\% lower, order rework proportion is 45.52\% lower, and the calendar-based throughput proxy is 71.73\% higher. The code-level reproduction confirms these numbers; it does not establish an analytical or industrial mechanism that would generate the assumed factor settings.

Consequently, the scenario comparison demonstrates how the specified sampler responds to assumptions about redesign. It cannot quantify the benefit of adding PM, XAI, or CNA. Such attribution would require interventions generated through reproducible analytical rules and comparison under common operational constraints, or deployment evidence.

\subsection{Statistical Validation of the Simulation Results}
\label{subsec:statistical_results}

Re-execution of the original scenario sampler reproduces all five headline percentages to the precision reported. Because scenarios reuse each replication's initial seed, the audit applies the blocked procedure described in Section~\ref{subsec:statistical_validation}. Table~\ref{tab:paired_statistics} reports G5--G0 relative changes, paired-bootstrap intervals, and Holm-adjusted baseline-comparison p-values.

\begin{table*}[t]
\caption{Executed seed-blocked G5--G0 scenario comparison. Relative-change intervals use paired replication-block bootstrap resampling. P-values are adjusted across five baseline comparisons within each KPI. These are conditional model outcomes, not causal effects of analytical modules.}
\label{tab:paired_statistics}
\centering\small
\begin{tabular}{lrrr}
\toprule
Indicator & Change (\%) & Bootstrap 95\% interval (\%) & Adjusted p\\
\midrule
Lead time reduction & 32.40 & [32.26, 32.54] & $ 9.31 \times 10^{-9} $ \\
Workload-proxy reduction & 45.09 & [44.66, 45.53] & $ 9.31 \times 10^{-9} $ \\
Defect-rate reduction & 27.63 & [25.45, 29.83] & $ 8.63 \times 10^{-6} $ \\
Rework-rate reduction & 45.52 & [40.72, 50.00] & $ 8.64 \times 10^{-6} $ \\
Throughput-proxy increase & 71.73 & [71.72, 71.74] & $ 9.31 \times 10^{-9} $ \\
\bottomrule
\end{tabular}
\end{table*}

All five G5--G0 comparisons remain significant within this sampler. This establishes stochastic differences between scenarios with different assumed factors. It does not validate physical throughput or inventory, estimate the contribution of a software module, or establish that model diagnosis caused the factor settings. The audit therefore supports H1 only as a within-model scenario comparison; H2 requires interpretation against activity-only baselines, and the broader integrated deployment proposition remains to be established.

\section{Discussion}
\label{sec:discussion}

The contribution is a traceable connection between current-state reconstruction, structural and model-assisted diagnosis, candidate intervention definition, offline assessment, and future-state representation. Rather than treating the analytical technologies as independent endpoints, the workflow organizes their outputs around decisions about process redesign. Model explanations and centrality indicators identify associations and structurally important elements; they do not establish causal intervention effects on their own.

The operational and analytical results have distinct scopes. Code inspection shows that the scenario sampler does not load the fitted ML models, explanations, or network metrics to determine its factors. Consequently, the six named configurations are assumed intervention scenarios rather than an executed component-ablation experiment. The complete G0--G5 comparison concerns a synthetic manufacturing simulation. The reductions of 32.40\% in lead time, 45.09\% in WIP, 27.63\% in defect rate, and 45.52\% in rework rate, with a throughput increase of 71.73\%, are simulated estimates conditional on the modeled system and intervention settings. The within-model statistical comparison addresses stochastic replication variability. It does not establish that these exact magnitudes transfer to physical production or that the simulator reproduces every industrial exception.

The numerical design compares predefined intervention factors associated conceptually with progressively enriched workflows. It does not isolate causal effects of individual modules, nor show that adding software alone improves production. A factorial design, component-removal comparison, and intervention-cost analysis would help separate analytical benefits from differences in the operational changes selected under each configuration.

The BPI benchmarks extend the analytical assessment to financial administration and procurement. An active-case count can provide an administrative workload analogue to WIP, but it is not physical inventory. Similarly, repeated administrative activities are repetition proxies unless corrective intent is confirmed, and measured processing time is not automatically value added. The domain-specific interpretations in Table~\ref{tab:domain_adaptation} are necessary for transfer of the workflow and prevent administrative classification scores from being presented as manufacturing validation.

The synthetic and real-log audits show that an activity-only rule suffices for perfect synthetic and BPI rework-label classification. Original BPI2017 results are reproduced, and the corrected reader retains valid timestamps previously dropped. The F1 = 1.0000 BPI2017 result is explained by the activity-defined target and a training-only activity signature. This control rules out the need for held-out graph construction to attain the perfect diagnostic result, but it does not establish a prospective forecasting capability. The near-perfect DQI operational-proxy classification has a related activity-level construct limitation. The results therefore support discussion of retrospective diagnostic separability rather than a claim of perfect prediction of future rework or physical bottlenecks. Prospective validation would require explicit decision timestamps and outcome horizons, prefix-only features, causal graph construction, and case-boundary handling.

\subsection{Resilience and Sustainability Implications}
\label{subsec:resilience_sustainability}

Industrial resilience concerns the capacity to sustain and recover operations under disruption. The workflow could support this objective by evaluating resource failures, demand surges, routing exceptions, and incomplete records as explicit intervention scenarios. Relevant outcomes would include service-level degradation, backlog accumulation, recovery time, and performance over disruption severity. The present nominal scenario comparison does not measure recovery or demonstrate resilience gains; these are concrete directions for extending the evaluation.

Sustainability-integrated VSM provides a relevant framework for connecting operational process analysis with environmental information \cite{HorsthoferRauch2024SustainabilityVSM}. Reduced rework or waiting may motivate investigation of resource savings, but lower operational waste does not by itself demonstrate lower energy use or emissions. Increased throughput and added capacity can introduce trade-offs. A sustainability extension should therefore measure energy per conforming unit, material loss, water use where relevant, and emissions under stated system boundaries, together with worker workload and safety. These indicators would allow explicit comparison with cost, lead time, and quality rather than assuming that operational gains automatically imply environmental or social improvement.

\subsection{Deployment and Evidence Boundaries}

A manufacturing deployment should establish the process boundary, timestamp semantics, resource constraints, and independently observed quality criteria. Simulation parameters should be estimated from an initial operating period and checked against a separate period before candidate interventions are assessed. Physical implementation followed by longitudinal measurements would then evaluate whether modeled benefits persist under changing demand, product mix, equipment failure, and exceptional routing.

The present closed loop is analytical: diagnosis informs interventions, simulation estimates their consequences, and a revised VSD represents the selected configuration. Physical implementation and a subsequent operating period were not evaluated. The supplied scripts implement an activity DFG and classification comparisons; resource/value-flow networks, SHAP explanations, conformance checking, and a shared-resource simulation remain proposed capabilities. A synchronized DT architecture could extend the workflow \cite{GrievesVickers2017,Aheleroff2021DTaaS}, but live connectivity and recurrent automatic recalibration are not capabilities established by this study.

